Portfolio allocation is usually posed as a trade-off between expected growth and risk, following mean-variance
analysis~\cite{markowitz1952} and its dynamic extensions~\cite{merton1969}. Optimized weights amplify errors in the
estimated inputs~\cite{michaud1989}. Practical work therefore shrinks the covariance estimate~\cite{ledoitWolf2004}
or adds limits on weights and trading, and judges the result against simple rules such as equal weighting, which is
difficult to beat out of sample~\cite{demiguelGarlappiUppal2009}. The single-index model~\cite{sharpe1963} reduces
the inputs to each stock's intercept and market loading, and factor models add further common
factors~\cite{famaFrench1993}. The allocation used here keeps the single-index inputs for its risk limits and its
scores, but replaces the mean-variance objective with utility functions from consumer theory: the Cobb--Douglas
form~\cite{cobbDouglas1928}, the constant elasticity of substitution form~\cite{arrowEtAl1961}, applied once to
share counts and once to portfolio weights, and the Stone--Geary form with a minimum holding of each
stock~\cite{stone1954, geary1950}, four utilities in all. In each utility, a stock's preference weight is its score,
a number between $-1$ and $1$ equal to the hyperbolic tangent of the ratio of its expected growth rate from the
single-index fit to a penalty on its market loading whose strength depends on the recent direction of the market. Each
preference weight is thus a single number that states why a stock is held. Stocks with a negative score receive no
preference weight in the long-only utilities. The expected growth rate rewards a stock's trend, and trend-following
signals in prices have a long empirical record~\cite{brockLakonishokLeBaron1992, moskowitzOoiPedersen2012}. The
penalty on large loadings, which is harsher when the market is falling, follows the evidence that portfolios tilted
away from high-beta stocks have earned higher risk-adjusted performance than portfolios tilted toward
them~\cite{frazziniPedersen2014}. Three benchmarks run under the same limits: the minimum-variance portfolio, the
tangency portfolio and a $130/30$ long-short portfolio, which holds $130\%$ of wealth long and $30\%$ short and so
relaxes the long-only constraint of the utility portfolios~\cite{jacobsLevyStarer1998}. Choosing which stocks to
hold before weighting them is a separate decision, and one of the selection rules studied here ranks the stocks of
each sector with a multi-armed bandit~\cite{lattimoreSzepesvari2020}, whose arms are the sector's stocks and whose
rewards are their growth in excess of the market over past 21-day windows.

With proportional trading costs, the optimal rebalancing policy leaves the holdings alone inside a no-trade region
and trades only at its boundary~\cite{magillConstantinides1976, davisNorman1990}. When expected growth rates are predictable, the
optimal policy instead trades a fraction of the gap between the current portfolio and the portfolio the investor
would hold without costs~\cite{garleanuPedersen2013}. Model-predictive control solves a finite-horizon optimization
from the measured state, applies the first decision, and solves again from the next measured
state~\cite{mayneEtAl2000}. Its portfolio applications include stochastic formulations with forecast models of asset
prices~\cite{herzogDondiGeering2007}, linear control policies for multi-period allocation~\cite{calafiore2008}, and
convex multi-period trading with costs and limits~\cite{boydEtAl2017}, all of which re-solve at every period.
Drawdown control has its own literature, in which exposure is adjusted continuously to keep wealth above a fixed
fraction of its running maximum~\cite{grossmanZhou1993}. The controller used here solves a single-period allocation
at each re-plan, forecasts the value of the resulting holdings over a $21$-day horizon, and keeps the holdings
unchanged until a trigger fires. Its default schedule is the horizon itself; a re-plan comes earlier when the
measured portfolio value leaves a band around the forecast, so that a forecast miss brings trading forward, or when
the drawdown reaches a limit. A re-plan re-solves the allocation from the current holdings and starts a new
forecast; it does not by itself reduce exposure.

Evaluating trading rules on the historical record carries a selection risk: when many rules or settings are tried on
the same path, the best measured result overstates the result to expect~\cite{baileyEtAl2014, harveyLiu2015}.
Synthetic data can test rules on paths that were not observed~\cite{assefaEtAl2020, jordonEtAl2022}, and the test is
informative to the extent that the generator reproduces the stylized facts of daily growth rates, in particular
heavy tails and volatility clustering~\cite{mandelbrot1963, cont2001}. Block bootstrap resampling keeps short-range
dependence~\cite{politisRomano1994}. Conditional variance models such as GARCH reproduce volatility
clustering~\cite{engle1982, bollerslev1986}. Regime-switching models, including hidden Markov models, represent
shifts between market states~\cite{hamilton1989}, and neural generators learn the joint distribution
directly~\cite{wieseEtAl2020}. Reviews of time-series synthesis note that a realistic distribution of daily values
does not imply realistic behavior over time~\cite{stengerEtAl2024}. Earlier work built the generator used here, a
hidden Markov model for each asset whose hidden states are bins of the daily growth rate, with heavy-tailed draws in
each state and occasional episodes in which the chain is held in its most extreme states for a random number of
days, producing bursts of large moves~\cite{alswaidan2026variance}. The same work combined the draws of many assets
through a common market factor, rescaling each asset's own draw so that adding the market term does not inflate the
asset's total variance. The present study adds drifting market loadings and sector factors to that combination, and
uses the result only as a market for the complete allocation and rebalancing system to trade in; the controller's
own forecast remains a geometric Brownian motion. The generator reproduces heavy tails more closely than volatility
clustering, and the evaluation carries that limitation.


